% Expository adaptation, October 7, 2026. Original mathematics: OpenAI, September 23, 2026.
% Reorganized and expanded for pedagogy; see README.md for the pinned source and scope.
\section{Detailed arithmetic constructions}

Section~\ref{sec:arith} introduced the signed-digit representation and stated the numerical guarantee used by the controller. This appendix proves the primitive operations from which that guarantee is assembled. For each operation, there are two separate questions: how accurately a finite computation approximates its exact value, and how much local information survives a request to one of its operands. We follow the same order for the elementary operations, multiplication, reciprocals, and indexed reductions.

\subsection{Compare finite truncations before taking a digit}

We first implement the elementary operations on finite truncations, and then bound their distance from the exact answer. Comparisons are between finite rational values; the algorithm never has to decide equality of two arbitrary exact reals. The truncation margin absorbs any change in which operand a maximum or minimum selects.


\begin{lemma}[Finite comparisons and bounded elementary operations]
\label{lem:arith-elementary}
The normalized operations of taking a pair average, a maximum or
minimum, and a fixed binary shift admit deterministic integers
satisfying Equation~\eqref{eq:arith-approximation}.  Their residue
procedures, for any fixed $q$, have the following properties.  There
are constants $C_0,C_1$ such that an operand place $j$ obeys
$j\leq t+C_0$, its suspended local record has length at most
$C_0+C_1(t-j)$, and the intrinsic local space is $O(t+1)$.
Fixed finite compositions of these normalized operations have the
same bounds, with enlarged constants.  Fixed Boolean operations on
returned digit codes, including sign decoding and conditional selection
of code bits, add only constant local overhead.
\end{lemma}

\begin{proof}
First consider comparing two truncations through $K=t+h$, where $h$
is a sufficiently large fixed integer.  Their difference digits have
magnitude at most 12.  Scanning from place zero to $K$, update a signed
prefix by $p\leftarrow2p+\delta$.  While $|p|\leq12$, retain $p$
exactly.  Once $p>12$ or $p<-12$, retain only the corresponding
absorbing sign state.  This sign is correct for every finite
continuation: the remaining suffix, in the current scale, has
magnitude strictly less than 12.  The state space is finite, so each
transition has a fixed bounded-fan-in Boolean circuit.  The sign
at the end determines the exact order of the two finite truncations.
The path from that sign bit to an input at place $j$ has length
$O(K-j+1)=O(t-j+1)$.

Let $v_t$ be the average, maximum, or minimum of the two
truncations, and put $N_t=\lfloor2^tv_t\rfloor$.  Each operation is
1-Lipschitz for the sup norm on its operands.  Consequently
$|v_t-v|\leq6\,2^{-K}$, where $v$ is the corresponding exact output.
Choosing $h$ so that $6\,2^{-h}<1/2$ gives
$|N_t-2^tv|<3/2$.

To obtain $N_t\bmod2^q$, only a fixed number of low places of the
selected truncation are needed.  More explicitly, at scale $2^K$,
terms of input place $j\leq t-q$ are multiples of $2^{q+h}$.
Modulo this modulus the remaining $q+h$ places suffice.  For an
average with integer input prefixes $A_K,B_K$ at scale $2^K$,
use $N_t=\lfloor(A_K+B_K)/2^{h+1}\rfloor$.  Its residue needs
modulus $2^{q+h+1}$ and one additional retained place:
only places $j\leq t-q-1$ can be omitted.  The finite number
of remaining digits can be added and shifted by a constant-size
Boolean circuit.  A maximum or minimum first uses the comparison
bit to select these residue bits.  Its long comparison branch has
the path bound proved above, and its residue branch reaches only
places within a fixed distance of $t$.  A fixed shift changes these
places by a fixed amount.  Fixed normalization factors are handled
in this way whenever an intermediate range exceeds $[-1,1]$.

The topology of the comparison chain and the low-residue circuits
depends only on the root clock, the fixed margins, and the gate
path.  Regenerating a place by replaying that path needs no saved
binary place counter at each transition.  Lemma~\ref{lem:arith-gate-evaluation}
therefore gives the suspended-space bounds.  Their longest paths
have length $O(t+1)$.

For composition, use one common coefficient of $(t-j)$ for the
finite primitive list.  Increasing such a coefficient by $\Delta$
requires increasing its constant term by at most $\Delta C_0$,
because the operand overshoot is at most $C_0$.  With these common
coefficients, the clock differences telescope along the composed
path.  The finite number of constant terms remains a constant.
\end{proof}

We will also use the following explicit Boolean arithmetic facts.
A carry bit in binary addition is described by a generate/propagate
pair $(g,p)$.  Composition of adjacent intervals is the associative
operation
\[
 (g_2,p_2)\circ(g_1,p_1)
   =(g_2\lor(p_2\land g_1),\ p_2\land p_1).
\]
Balanced trees of these pairs give every carry in $O(\log(W+1))$
depth for $W$-bit words.  A balanced tree adding $R$ such words modulo
$2^W$ has depth
\begin{equation}
 O\bigl(\log(R+1)\log(W+1)\bigr).
 \label{eq:arith-adder-depth}
\end{equation}
The trees can be padded by zero words to powers of two, so their
gate addresses are computable by ordinary index arithmetic.
For counts of $R$ Boolean bits one can do better: constant-depth
full adders replace three same-width rows by two rows of their sum,
with carries placed in the next column.  In $O(\log(R+1))$ rounds
only two rows remain.  The word width is $O(\log(R+1))$; a final
carry computation gives total depth $O(\log(R+1))$.  These
constructions use bounded fan-in throughout.

\subsection{Multiply without retaining a long partial product}

Expanding a product of digit streams produces many pairs of input digits. A straightforward sum can leave a long partial result stored while requesting another digit. We instead partition the pairs according to the smaller of their two indices. In a group whose smaller index is large, both digits are early enough relative to the requested output place to pay for the circuit path leading to them.

The proof has four parts: define finite band sums, bound the omitted tails and rounding error, compute only the residues needed for the output digit, and charge each circuit path to the difference between its output and input places.

\begin{lemma}[Multiplication]
\label{lem:arith-multiplication}
For normalized exact input streams $x$ and $y$, multiplication has all
the conclusions of Lemma~\ref{lem:arith-elementary}.  In particular,
its suspended storage at an operand place $j$ is
$O(1)+c(t-j)$, despite the potentially long product expansion.
\end{lemma}

\begin{proof}
First partition the formal product and specify which pairs are retained. Fix $q$ and write $A=6$.  Partition ordered pairs $(i,j)$ of
nonnegative places into a zero band, where $\min(i,j)=0$, and
positive bands
\[
 \mathcal P_D=\{(i,j):D\leq\min(i,j)<2D\},
 \qquad D=1,2,4,\ldots.
\]
Set $h_0=h_\circ$ for the zero band and
$h_D=h_\circ+\lfloor D/4\rfloor$ for positive bands.  The fixed
integer $h_\circ$ will be chosen below.  Retain in band $D$ only
pairs with $i+j\leq t+h_D$, and let
\begin{equation}
 S_D=\sum_{\substack{(i,j)\in\mathcal P_D\\i+j\leq t+h_D}}
           x_i y_j2^{t+h_D-i-j}.
 \label{eq:arith-band-integer}
\end{equation}
For the zero band the same formula uses its defining pair set.
These are mathematical integers; no bound on their full binary
length is assumed in the residue computation.

The omitted tails can be bounded without forming the band sums. At any fixed diagonal $i+j=s$, a positive band contains at most
$2D$ ordered pairs.  Indeed the smaller place has $D$ possibilities,
and an off-diagonal pair has two orientations.  The zero band has
at most two pairs.  Thus, in units of $2^{-t}$, the total omitted
product terms have absolute sum at most
\begin{equation}
 E_{\mathrm{omit}}
 \leq 2A^2\sum_{D\in\{0,1,2,4,\ldots\}}(D+1)2^{-h_D}.
 \label{eq:arith-band-tail-error}
\end{equation}
The geometric summation of the diagonal weights gives this bound
also when the truncation threshold lies below the first possible
diagonal.  Both the series in Equation~\eqref{eq:arith-band-tail-error}
and $\sum_D2^{-h_D}$ converge.  Choose $h_\circ$ sufficiently large
that their displayed combined error bound, including coefficient
$2A^2$, is less than $1/2$.

There are only finitely many nonempty retained bands.  For a
positive one, $2D\leq t+h_D$, whence
\begin{equation}
 D\leq\frac47(t+h_\circ).
 \label{eq:arith-last-band}
\end{equation}
The left side of the nonemptiness condition grows monotonically
faster than its right side along the dyadic bands.  We may therefore
stop at the last band satisfying this inequality, retaining empty
bands within the finite prefix if necessary.

Starting with the most distant retained band, define integers
recursively by
\begin{equation}
 F_D=S_D+
       \left\lfloor 2^{h_D-h_{D^+}}F_{D^+}\right\rfloor,
 \qquad
 N_t=\left\lfloor2^{-h_0}F_0\right\rfloor,
 \label{eq:arith-band-recursion}
\end{equation}
where $D^+$ is the next band, and the last missing tail is zero.
One internal floor, divided back to scale $2^t$, changes the sum
by less than $2^{-h_D}$.  Induction down the bands consequently
shows that the difference between $2^{-h_0}F_0$ and the sum of
all retained terms in scale-$t$ units is at most $\sum_D2^{-h_D}$.
Equation~\eqref{eq:arith-band-tail-error} and the choice of
$h_\circ$ make the error before the final floor less than $1/2$.
The final floor costs less than one more unit.  Therefore
Equation~\eqref{eq:arith-approximation} holds for $z=xy$.

The integers in this recurrence define the approximation; they need not be stored in full. We next implement Equation~\eqref{eq:arith-band-recursion} using residues.  At band $D$ use modulus $2^{q+h_D}$.  An integer term
in Equation~\eqref{eq:arith-band-integer} with $i+j\leq t-q$
is a multiple of this modulus and can be omitted from the circuit.
It remains present in the mathematical integer $S_D$.
Only the diagonals
\begin{equation}
 t-q<i+j\leq t+h_D
 \label{eq:arith-relevant-diagonals}
\end{equation}
need to be processed, giving
$O((D+1)(h_D+q+1))=O((D+1)^2)$ candidate pairs.

Signed floor shifts are well defined on the required residues.
For $h'\geq h$ and $a-b=k2^{q+h'}$,
\begin{equation}
 \left\lfloor\frac{a}{2^{h'-h}}\right\rfloor
 -\left\lfloor\frac{b}{2^{h'-h}}\right\rfloor
       =k2^{q+h}.
 \label{eq:arith-signed-residue-floor}
\end{equation}
This identity holds for all signed integers $a,b$.  Thus a next-band
residue modulo $2^{q+h'}$ supplies the correctly floor-shifted
tail modulo $2^{q+h}$.  Apply the same observation to the final
shift by $h_0$ to obtain $N_t\bmod2^q$.
In the finite binary circuit, this shift just discards the lowest
$h'-h$ bits of the canonical $(q+h')$-bit residue.  The displayed
identity proves that this is valid also when the underlying integer
is negative.

For completeness, the finite gate structure can be indexed without
storing an absolute operand place at each gate.  For a positive
band enumerate triples $(r,s,o)$, where
$D\leq r<2D$, $t-q<s\leq t+h_D$, and $o\in\{0,1\}$.
The triple supplies $(r,s-r)$ if $s-r\geq r$ and $o=0$;
it supplies $(s-r,r)$ if $s-r>r$ and $o=1$; otherwise it supplies
a zero term.  This enumerates the relevant pairs exactly once.
For the zero band use $r=0$ with the same orientation rule.
Negative candidate places are invalid and supply zero.
The offsets $r-D$ and $s-(t-q+1)$ have $O(\log(D+2))$ bits.
They and $o$ are obtained from paths in padded binary selection trees.
The full places are then computed from the root clock in scratch.

Each nonzero term is a constant-bit signed product $x_i y_j$,
shifted into a word of width $W_D=q+h_D=O(D+1)$ and interpreted
modulo $2^{W_D}$.  A bit of this word is a fixed Boolean function
of the two bounded digit codes, the bit index and the known shift;
negative products use their ordinary residue, equivalently sign
extension before reduction.  Add these words and the shifted tail
using Equation~\eqref{eq:arith-adder-depth}.  A band contributes
depth $O(\log^2(D+2))$.  Bit selection for the shifted tail is
wiring and does not require evaluating a full tail word.

It remains to connect circuit depth with the digit place of an operand. This is the reason for grouping by the smaller input index. The sum of these depths through the bands up to $D$ is
$O(\log^3(D+2))$, hence at most $C(D+1)$ for a fixed $C$.
For an original input digit in a positive band, its partner place
is at least $D$.  Equation~\eqref{eq:arith-relevant-diagonals} gives
\begin{equation}
 j\leq t+h_D-D
       \leq t+h_\circ-\frac34D.
 \label{eq:arith-multiplication-latency}
\end{equation}
Equivalently, $t-j\geq\tfrac34D-h_\circ$, so the full path cost
through these bands satisfies
\[
 C(D+1)\leq C\left(1+\frac43h_\circ\right)
                   +\frac{4C}{3}(t-j).
\]
This gives $C_0+C_1(t-j)$ with fixed constants.  The zero band
has constant depth and only constant overshoot.  All its gate
descriptions, term coordinates, bit positions and tail shifts are
computable by replaying paths and doing integer index arithmetic
from $t$ and $q$.  Lemma~\ref{lem:arith-gate-evaluation} converts these
depth bounds to suspended storage.  Finally,
Equation~\eqref{eq:arith-last-band} bounds every intrinsic path
and local word width by $O(t+1)$.
\end{proof}

\subsection{Take a reciprocal using an early prefix and a small remainder}

For a reciprocal, a prefix through roughly half the requested digit place is enough to define a linear approximation whose remaining error is quadratically small. Forming that prefix requires a long word. The memory bound is still possible because the word is present only while the algorithm queries early input places, where the output-to-input digit difference is large enough to pay for it.

The crucial order is to compute a needed coefficient digit from the early prefix, release the prefix word, and only then request a late digit of the input remainder. Treating the prefix as a permanent local constant would lose the storage bound.

\begin{lemma}[Reciprocal on a positive compact interval]
\label{lem:arith-reciprocal}
Fix constants $0<a\leq b$.  On inputs $x\in[a,b]$, a reciprocal,
after a fixed known normalization, has the same operand-clock,
suspended-record, and intrinsic-space bounds as in
Lemma~\ref{lem:arith-elementary}.  Its local record may contain an
$O(t+1)$-bit word only while querying input places
$j\leq\lfloor t/2\rfloor+O(1)$.
\end{lemma}

\begin{proof}
The input can first be normalized by a fixed binary scale.  It is
enough to prove the result with a fixed digit alphabet and fixed
positive interval; all constants below may depend on $a,b$.
Choose a fixed power of two $C$ large enough that the desired output
$z=1/(Cx)$ lies in $[0,1]$ and that the coefficients below have
absolute value at most one.  Put
\[
 k=\lfloor t/2\rfloor+c,\qquad x^0=x^{[k]},\qquad e=x-x^0,
\]
where $c$ is a sufficiently large fixed integer.  By
Equation~\eqref{eq:arith-tail}, $|e|\leq 6\,2^{-k}$.  Choosing $c$
first ensures $x^0\in[a/2,2b]$ and $|e|\leq1$ at every $t$.
Increasing $C$ if necessary gives bounded coefficients
\[
 \alpha=\frac1{Cx^0},\qquad \beta=-\frac1{C(x^0)^2}.
\]
The exact algebraic identity
\begin{equation}
 \frac1{Cx}=\alpha+\beta e+
               \frac{e^2}{C(x^0)^2x}
 \label{eq:arith-reciprocal-taylor}
\end{equation}
shows that the linear expression $T_t=\alpha+\beta e$ differs
from $z$ by at most a fixed constant times $2^{-2k}$.
Choose $c$ large enough that
\begin{equation}
 2^t|T_t-z|\leq\frac14
 \quad\hbox{for every }t\geq0.
 \label{eq:arith-reciprocal-remainder}
\end{equation}
Clipping $T_t$ to $[-1,1]$ cannot increase this error; denote the
clipped expression by $\widehat T_t$.

For this fixed original $t$, the coefficients $\alpha,\beta$ are
exact rational numbers.  Their digits are produced by integer
arithmetic, as described below.  The suffix $e$ has the exact stream
with digit zero through $k$ and digit $x_j$ thereafter.  Thus
Lemmas~\ref{lem:arith-elementary} and
\ref{lem:arith-multiplication} provide exact streams for
$\widehat T_t$, using only a fixed number of normalized operations.
Explicitly, form the product $\beta e$, then the pair average
$u=(\alpha+\beta e)/2$, clip $u$ to $[-1/2,1/2]$ by a maximum
and a minimum, and finally apply the fixed shift by 2.  Every
exported intermediate value is normalized, and the final value is
$\widehat T_t$.  Forming this expression introduces no additional
reciprocal operation.

Let $P_s(\widehat T_t)$ be the integer approximation furnished by
those operations at place $s=t+3$, and define
\begin{equation}
 N_t(z)=\left\lfloor P_{t+3}(\widehat T_t)/8\right\rfloor.
 \label{eq:arith-reciprocal-approximation}
\end{equation}
The cutoff $k$ and the coefficients remain those of the original $t$,
not those of $s$.  Equations~\eqref{eq:arith-approximation} and
\eqref{eq:arith-reciprocal-remainder} give
\[
 |N_t(z)-2^tz|<1+\frac28+\frac14=\frac32<2.
\]
To obtain this integer modulo $2^q$, ask for the approximation
$P_{t+3}$ modulo $2^{q+3}$ and use
Equation~\eqref{eq:arith-signed-residue-floor}.

The approximation error is now controlled. We separately verify the lifetime of the operand-dependent words, since their rational form alone does not make them public constants. It remains to justify storage in the coefficient computations.
The signed numerator of $x^0$ in units $2^{-k}$ has $O(k+1)$
bits.  Form it by querying places $0,\ldots,k$ and updating the
integer prefix.  The numerator, a place counter, and all data that
must survive the next query are part of the local record.
To return a coefficient digit of place $r$, numerator squaring and
integer division require $O(k+r+1)$ bits.  Here $r\leq t+O(1)$,
because the enclosing expression has fixed depth and constant
overshoot.  After the input prefix is obtained, these integer
computations make no further operand queries.  Standard binary
long multiplication and division, performed by bounded loops with
their partial products or remainder, use only this linear number
of bits.  For $r>0$, computing floors of the coefficient scaled at
$r$ and $r-1$ gives its exact difference digit; at $r=0$ use the
place-zero floor alone, as in Lemma~\ref{lem:arith-residue-stream}.
Clipping is unnecessary
because the coefficients were normalized in advance.

Every original-input query during this formation has
$j\leq k\leq t/2+O(1)$.  Hence
\begin{equation}
 t+1\leq C_2+C_3(t-j)
 \label{eq:arith-early-word-payment}
\end{equation}
for fixed constants, at all such queries.  The enclosing gate path
also has length $O(t+1)$, so the complete suspended local record
satisfies the required bound there.  Equivalently, the coefficient
branch is charged against the original reciprocal clock $t$, not
against its possibly small coefficient digit place $r$.

On a suffix branch with $j>k$, the coefficient word is absent.
The Boolean evaluation asks for a coefficient digit at a separate
leaf, returns its bounded code bits, and discards the coefficient
word.  If the other child of a gate subsequently asks for a late
suffix digit, only the bounded pending bits and gate choices are
retained.  The path bound from the preceding arithmetic lemmas
therefore gives $O(1)+c(t-j)$ also on this branch.  Repeated
coefficient requests rebuild the word.  The parameter $k$ and the
choice of coefficient are regenerated from the original $t$ and
the saved gate path; no earlier cursor address is required.
All intrinsic records have $O(t+1)$ bits.
\end{proof}

The coefficients in this proof are operand-dependent, even though they
are rational once $x^0$ is known.  They therefore do not qualify as
public coefficient leaves.  A coefficient request first forms its
charged early prefix, then uses scratch for integer arithmetic, returns
only the requested bounded digit, and releases that prefix before a
suffix branch is evaluated.  The original reciprocal clock $t$ remains
regenerable throughout.  This order of operations is the lifetime rule
used in the suspended-record bound.  In particular, $j\le\lfloor t/2\rfloor+c$
gives the explicit payment $t+1\le2(t-j)+2c+1$ for the early word.

\subsection{Order statistics, sums, and finite-accuracy access tests}

The remaining operations have an explicitly indexed domain of size at most exponential in the local descriptor length. We can recompute a value whenever it is compared or added, rather than store the whole domain. Numerical order statistics may select different approximating entries at different digit places, but they must approximate one fixed exact order statistic. By contrast, a structural access test has a fixed precision chosen once for that test.

\begin{lemma}[Indexed finite ranks]
\label{lem:arith-rank}
Fix $\kappa>0$ and $d\geq1$.  Consider at most $2^{\kappa d}$
normalized values, indexed by descriptors of length at most
$\kappa d$, and an order statistic specified by an $O(d)$-bit
index.  The semantic list length $n$ is public; the requested index
is in $\{1,\ldots,n\}$, or an out-of-range request returns a
prescribed normalized constant.  Absent descriptors may be assigned
prescribed constant dummy values by structural validity predicates.
This order statistic has deterministic integer approximations $N_t$
satisfying Equation~\eqref{eq:arith-approximation}.  Their residue
procedures have constants $a,b,c$, depending only on the descriptor
bounds, such that an operand digit query at $j$ obeys
\begin{equation}
 j\leq t+cd,\qquad
 L\leq ad+b(t-j),
 \label{eq:arith-local-bound}
\end{equation}
where $L$ is the total suspended local record length.  A structural
validity query is assigned clock zero and has suspended length
$O(d+t)$.  The intrinsic local length is $O(d+t)$.
\end{lemma}

\begin{proof}
For $1\le k\le n$, the $k$th largest of the semantic $n$-entry list
is its $(n-k+1)$st ascending order statistic, before any extra padding.
Pad to a fixed descriptor domain whose cardinality is a power of
two, adding entries of value $+1$.  These entries do not change
any of the first $n$ ascending order statistics.  An empty list or
an out-of-range request returns its specified default immediately.
Replace invalid original descriptors by the prescribed dummies.
An operand access evaluates validity first and enters the actual
value branch only when valid; an invalid access returns the dummy's
digit.  Thus a dummy does not require a query to an undefined value.
For each valid numerical value use its truncation through
$K=t+h$ for fixed large $h$.  Order these finite rational values,
breaking ties by descriptor index.  Compare every pair using the
signed comparison of Lemma~\ref{lem:arith-elementary}.  For each
descriptor count how many descriptors precede it, then compare
that count to the requested rank minus one.  Exactly one
descriptor is selected.  Select its low-residue bits using Boolean
AND and OR trees and finally floor to output scale $2^t$.
The selected descriptor is internal to this residue computation.  It
may change with $t$; the exact order statistic being approximated is
fixed by the semantic list and the requested index.

If all coordinates of two finite lists differ by at most $\eta$,
their corresponding order statistics differ by at most $\eta$.
Indeed, at least the specified number of the perturbed entries
are no larger than the original order statistic plus $\eta$;
interchanging the lists gives the reverse inequality.  Here
$\eta=6\,2^{-K}$.  Thus the selected truncation, followed by the
final floor, satisfies Equation~\eqref{eq:arith-approximation}
for sufficiently large fixed $h$.  Descriptor ties affect neither
this estimate nor the uniqueness of the finite selection.
This argument perturbs the entire padded exact list.  Some truncated
original entries may exceed $+1$ slightly, but their error, and hence
the error of the selected padded order statistic, is still at most
$\eta$.  The padding argument requires no separate claim that every
finite truncation remains in $[-1,1]$.

Each count has $O(d)$ bits.  Carry-save counting as described
above, the comparison to the requested rank, and selection among
$2^{O(d)}$ descriptors add $O(d)$ Boolean depth above the
pairwise comparison bits and selected residue bits.  A
place-$j$ input to a pairwise comparison consequently has path
length $O(d)+O(t-j)$; a residue input is within a fixed distance
of $t$ and has path length $O(d)$.  The paths to validity predicates
have length at most the whole circuit depth $O(d+t)$, and the
predicates themselves are separate structural leaves of clock zero.

Each gate is specified by its pair of descriptor indices, its
counting-tree or selection-tree position, and a comparator or bit
position.  The paths determine these indices in $O(d+t)$ scratch,
using only the fixed domain and root parameters.  No numerical
comparison is needed to determine the circuit topology.  Apply
Lemma~\ref{lem:arith-gate-evaluation}.  The constants in
Equation~\eqref{eq:arith-local-bound} can absorb the fixed digit
overshoot because $d\geq1$.
\end{proof}

For the compression ranks, copies of a previously computed baseline
are ordinary indexed operands: their digit interfaces query the
baseline's exact stream one bounded digit at a time.  A slot that
chooses between an entry and the baseline first evaluates its
structural flag and then queries the selected, already-defined stream.
Treating this as a defined operand interface leaves no absent value
for the rank procedure to evaluate.
The required baseline copies belong to the semantic list; the extra
power-of-two padding in the proof is then supplied by the constant
$+1$.  The structural choice adds only a fixed conditional branch,
so the same digit-clock and suspended-record bounds apply.

A rank over the full incoming vertex universe need not have a $2^{O(d)}$ descriptor domain. That operation uses the terminal-program construction of Lemma~\ref{lem:ctrl-score-rank}, rather than the finite-domain result proved here.

A sum is evaluated by a balanced tree of normalized pair averages, followed by a known scale adjustment. This avoids retaining a sequential accumulator across late operand queries. The number of tree levels is paid for by the descriptor length.

\begin{lemma}[Bounded sums and known scales]
\label{lem:arith-sums-scales}
Suppose an immediate descriptor domain has at most $2^{\kappa d}$
terms, all required binary scale exponents have magnitude at most
$\kappa d$, and every intermediate output has a known range bound
$2^{\kappa d}$.  A sum over an empty domain is zero, and an average
is taken over a public nonempty domain.  Balanced sums or averages of these terms, followed
by the required known rescalings and bounded clipping, satisfy
Equation~\eqref{eq:arith-local-bound} and have intrinsic local
space $O(d+t)$.
\end{lemma}

\begin{proof}
Choose one common known power-of-two range bound for the immediate
terms and divide each term by it, so every term lies in $[-1,1]$.
Pad their number to a power of two $2^r$, $r=O(d)$, and use a complete
binary tree of pair averages.  Every internal value stays normalized.
Along one path there are $r$ elementary operations.  Choose common
coefficients in their bounds, as in the last paragraph of the
proof of Lemma~\ref{lem:arith-elementary}.  The constant terms
sum to $O(r)$ and the differences of successive digit clocks
telescope.  Thus the local length is $O(d)+b(t-j)$ and the
total overshoot is $O(d)$.  Restore the common normalization scale
only after the tree of pair averages.  For a sum, combine this factor
with the known padding count $2^r$ in one binary shift of magnitude
$O(d)$.  For an average over an original domain of size $m$, combine
the scale restoration with multiplication of the padded average by
the public factor $2^r/m$.  With minimal padding the latter factor
is in $[1,2)$, so it requires only one normalized public multiplication
and a fixed shift in addition to the common scale restoration.

More explicitly, for a shift by $2^v$, truncate the operand at
$K=\max(0,t+v+h)$.  If $t+v+h\geq0$, its scaled truncation
error in scale-$t$ units is at most $6\,2^{-h}$.  If the scaled
output is known to be normalized, floor in the desired output
scale.  Otherwise, when a bounded clip is required, apply that
clip to the scaled truncation before flooring.  For the
normalized clip to $[-1,1]$, for example, define the approximation
using $\min(1,\max(-1,2^v x^{[K]}))$; its error cannot increase.
This shift-and-clip procedure does not first export an unbounded
intermediate as a normalized stream.
The low residue requires the constant number of places near
$t+v$, and a comparison needed for clipping has a constant-state
chain through $K$.  Comparisons to the scaled clipping endpoints
are equivalently comparisons of $x^{[K]}$ to the known endpoints
divided by $2^v$; their finite constants and shifted place data
are public.  When $v\geq0$ these endpoints belong to $[-1,1]$;
when $v<0$ the normalized operand is already in the clip interval
after scaling, so no clipping comparison is needed.  Its path cost is
$O(1)+O(\max(0,t+v-j))$, which is at most
$O(d)+O(t-j)$ on an actual operand query.  If $t+v+h<0$, zero
is already a sufficiently accurate approximation after increasing
$h$, and no operand query is needed.  A fixed number of such
rescalings and clips preserves the bounds by telescoping.
Known non-power-of-two coefficients can instead be handled as
public normalized streams and bounded multiplication.
\end{proof}

An access decision does not need to distinguish values exactly at one sharp threshold. The allowed gap between acceptance above $2u$ and rejection at or below $u$ makes a fixed finite precision sufficient. This is the numerical basis for the support margins used in the table and compression constructions.

\begin{lemma}[Finite buffered tests]
\label{lem:arith-buffered-test}
Let $d\geq1$, let $s$ and $u$ have known magnitude bound $2^{\kappa d}$,
and suppose $u\geq\rho\geq2^{-\kappa d}$.  There is a deterministic
finite test that accepts every $s>2u$ and accepts only $s>u$.
It uses $O(d)$-place digit queries after normalization, with
intrinsic and suspended local lengths $O(d)$.  Its result is
independent of any subsequent numerical output clock.
\end{lemma}

\begin{proof}
Compute finite rational approximations $\widetilde s,\widetilde u$
with individual errors at most $\rho/16$.  The known range and
lower bound on $\rho$ make $O(d)$ normalized digits sufficient.
Accept exactly when
$\widetilde s>\tfrac32\widetilde u$, using a fixed convention
at equality.  If $s>2u$, the left side minus the right side is
greater than $u/2-5\rho/32>0$.  If $s\leq u$, it is at most
$-u/2+5\rho/32<0$.  The finite comparison and all previous
arithmetic have the bounds already proved.  The test chooses its
precision from $d$ and $\rho$, not from an enclosing numerical
clock.  With its structural starting clock defined as zero, its
operand requests have $j=O(d)$; increasing $a$ in
Equation~\eqref{eq:arith-local-bound} pays also for the term $-bj$.
\end{proof}

